% Source: Wald GR, physical PDF p. 217
%         (printed p. 204).
% Coverage: A sequence approaching a Cauchy horizon.
% Figures/tables/footnotes: Figure 8.13; no tables or footnotes.
% Status: source-reviewed and compile-clean.
% Uncertainties: none.

\begingroup
\renewcommand{\thefigure}{8.13}
\begin{figure}[H]
  \centering
  \begin{tikzpicture}[x=.80cm,y=.80cm,>=Latex,line cap=round,font=\small]
    \draw[thick] (-1.55,.12) .. controls (-.55,.58) and (.48,.22) .. (1.25,.34)
      -- (3.05,1.95);
    \draw[thick] (1.25,.34) -- (.13,2.62) -- (3.05,1.95);
    \node at (1.55,1.29) {\(D^+(S)\)};
    \draw[->] (.76,2.90) .. controls (.63,2.52) and (.67,2.26) .. (.84,2.01);
    \draw[->] (1.90,2.91) .. controls (1.96,2.53) and (2.16,2.29) .. (2.39,2.13);
    \node at (1.30,2.97) {\(H^+(S)\)};
    \draw[->] (1.73,.06) -- (1.60,.37) node[below=.16cm] {\(S\)};
    \fill (-.22,1.76) circle (1.6pt) node[right] {\(p\)};
    \foreach \x/\y in {-.42/2.08,-.57/2.37,-.72/2.67,-.87/2.95} {
      \fill (\x,\y) circle (1.2pt);
    }
    \node at (-.61,3.29) {\(q_n\)};
    \draw[thick] (-1.70,-1.10) .. controls (-1.47,.44) and (-.98,1.51) .. (-.42,2.08);
    \draw[thick] (-1.53,-1.00) .. controls (-1.24,.64) and (-.89,1.89) .. (-.57,2.37);
    \draw[thick] (-1.34,-.89) .. controls (-1.05,.78) and (-.83,2.16) .. (-.72,2.67);
    \draw[thick] (-1.16,-.76) .. controls (-.92,.91) and (-.81,2.46) .. (-.87,2.95);
    \draw[->] (-1.08,1.56) -- (-1.59,1.31);
    \node at (-1.68,1.79) {\(\lambda_n\)};
  \end{tikzpicture}
  \caption{A point \(p\in H^+(S)\cap I^+(S)\) with a sequence of points
  \(q_n\in I^+(p)\) converging to \(p\) (see theorem~\ref{thm:8.3.5}).}
  \label{fig:8.13}
\end{figure}
\endgroup
